% !TeX root = ../教材题库_全部条目_转写工作稿.tex
% Source: XJ-B1-C05-TH-P156-Y0591
% 第5章 / 5.1 任意角和弧度制 / 教材 p.156、157 / PDF 第163、164页
% OCR未校对：公式与文字均可能有误；未完成数学TeX转写；禁止直接发给学生
% 原题及相邻原解截图（核对用）：../题目原图/XJ-B1-C05-TH-P156-Y0591-01.webp, ../题目原图/XJ-B1-C05-TH-P156-Y0591-02.webp
% 本文件只在转写完成后作为题干使用；OCR 条目尚不是可靠题干。
\begin{bankitem}{p.156 思考（含紧邻说明）}{第5章 / 5.1 任意角和弧度制 / 教材 p.156、157 / PDF 第163、164页}{XJ-B1-C05-TH-P156-Y0591}{OCR未校对：公式与文字均可能有误；未完成数学TeX转写；禁止直接发给学生}
\begin{ocrdraft}
[教材p.156 / PDF第163页]
\par
？思考万物皆变，万物皆动．有平动，有转动．平动量
\par
距离，转动量什么？以时钟为例，时钟现在的时间是10：10，经
\par
过40mrn后，时钟的分针转过了多少度？经过90min呢？
\par
我们知道，角可以看作是平面内一条射线绕着其端点从初始
\par
位置旋转到终止位置时所形成的图形．以前学习的角，其大小都
\par
在0。到360。之间，而在现实生活中，可以发现大量由旋转形成的角大于36的例
\par
如，行驶的自行车、汽车、火车的车轮，它们通常都要旋转很多圈．此外，平面上
\par
的旋转可以是顺时针方向，也可以是逆时针方向．因此，要准确地刻画旋转所成的
\par
角不仅要知道旋转的量，还要知道旋转的方向．
\par
为了表示不同旋转方向所形成的角，联想到正负数可表示具有相反意义的量，
\par
我们做如下规定：
\par
0
\par
一条射线绕着端点以逆时针方向旋转所成的角称为正角；
\par
“=420
\par
以顺时针方向旋转所成的角称为负角；不旋转所成的角称为
\par
0
\par
零角，用0。表示．零角的始边与终边重合．
\par
=一150。
\par
如图5．1一1中，角a是OA沿逆时针方向绕0点转动形
\par
0
\par
成的角，所以一420。；角是OA沿顺时针方向绕0点转动
\par
形成的角，所以一一150。
\par
这样，角的概念就推广到了任意角，包括正角、零角和负角．
\par
于是，上述“思考”中的答案分别是一240。，
\par
一540。．
\par
为了方便研究问题，常在平面直角坐标系内讨论角，为
\par
本书中提到的，若
\par
此取坐标原点为角的顶点，x轴的非负半轴为角的始边，那
\par
不特别说明，总是指角的
\par
么，角的终边落在第几象限，就说这个角是第几象限角；如
\par
顶点与原点重合，角的始
\par
果角的终边在坐标轴上，就认为这个角不属于任何一个象限．
\par
边与工轴的非负半轴重合．
\par
例如，如图5．1一2，45。角和405。角都是第一象限角，120。角
\par
[教材p.157 / PDF第164页]
\par
和一240。角都是第二象限角．
\par
120。
\par
45
\par
0
\par
405
\par
一240。
\par
（2）
\par
图5．1一2
\end{ocrdraft}
\end{bankitem}
